% =============================================================================
% rebuttal-template.tex — generic ONE-PAGE author-rebuttal scaffold
% Part of research-paper-skills / write-rebuttal. MIT licensed; original work,
% no venue style files copied.
%
% WHEN TO USE THIS FILE
%   - The venue mandates a one-page PDF rebuttal but ships NO official
%     rebuttal template, or you are drafting before porting.
%   - If the venue ships an official rebuttal kit (e.g. CVPR's author-kit
%     repository includes a rebuttal template), you MUST use that kit for the
%     actual submission — wrong template is a reject-without-review trigger.
%     Draft here if convenient, then port the body text into the official kit.
%
% HARD RULES THIS LAYOUT ENCODES (re-verify on the venue's author guidelines):
%   - exactly ONE page; over-length rebuttals are typically not reviewed
%   - two-column, 10pt, letterpaper — matches the common CV-venue look
%   - anonymous: no author names, affiliations, or identifying links
%   - no external links (code/video URLs are usually banned in rebuttals)
%   - do NOT touch \geometry, font sizes, or add negative \vspace to gain
%     space — template tampering is itself a desk-reject trigger
%
% BUILD + CHECK
%   pdflatex rebuttal.tex
%   python3 skills/write-rebuttal/scripts/check_budget.py pdf rebuttal.pdf --max-pages 1
% =============================================================================
\documentclass[10pt,twocolumn,letterpaper]{article}

\usepackage[letterpaper,margin=0.75in,top=0.9in,bottom=0.9in]{geometry}
\usepackage{times}            % Times text — the conventional rebuttal look
\usepackage[T1]{fontenc}
\usepackage{xcolor}
\usepackage{graphicx}
\usepackage{booktabs}         % only if you include a compact results table
\usepackage[pdfborder={0 0 0}]{hyperref} % internal refs only; NO external URLs

\pagestyle{empty}             % many venues require no page number on rebuttals
\setlength{\parindent}{0pt}
\setlength{\parskip}{3.5pt}

% --- rebuttal macros ---------------------------------------------------------
% \rpoint{R2.1}{abridged reviewer quote} — a reviewer point, bold ID + italic quote
\newcommand{\rpoint}[2]{%
  \vspace{2.5pt}\noindent\textbf{[#1]} \textit{``#2''}\par\nopagebreak}
% \resp{...} — your response paragraph (plain; keep one para per point)
\newcommand{\resp}[1]{\noindent #1\par}
% \rhead{Reviewer 1 (R1)} — per-reviewer block heading
\newcommand{\rhead}[1]{%
  \vspace{5pt}\noindent\textbf{\large #1}\par\nopagebreak\vspace{1.5pt}}
% Reviewer-colored inline IDs for cross-references like "see [R1.2]"
\newcommand{\rid}[1]{\textbf{[#1]}}

\begin{document}

% --- title line (NO author block — rebuttals stay anonymous) -----------------
\twocolumn[{\centering
  {\Large\bfseries Rebuttal for Submission \#XXXX:\\[2pt]
   Paper Title Goes Here\par}
  \vspace{8pt}
}]
\thispagestyle{empty}

% --- 2–4 line global summary --------------------------------------------------
We thank all reviewers for their careful reading. R1 and R3 find the method
novel and the evaluation thorough; the main concerns are a missing baseline
(\rid{R1.1}, \rid{R2.2}) and the scope of the ablation (\rid{R2.1}). We
address every point below, citing sections, tables, and line numbers of the
submitted PDF.

% --- shared concerns first (saves space when reviewers overlap) ---------------
\rhead{Shared concern: missing baseline B (\rid{R1.1}, \rid{R2.2})}
\resp{Table~2 (L412--418) already compares against B's strongest published
configuration; we agree the connection was easy to miss and will name it
explicitly in the caption. On the requested setting, the protocol of
Sec.~5.1 applies unchanged; we will add the full row in the revision.}

% --- per-reviewer blocks -------------------------------------------------------
\rhead{Reviewer 1 (R1)}

\rpoint{R1.2}{The complexity analysis in Sec.~4 seems to assume sparsity.}
\resp{Correct, and the assumption is stated at L231; Appendix~C gives the
dense-case bound. We will surface both in the main text (one sentence,
Sec.~4.3).}

\rhead{Reviewer 2 (R2)}

\rpoint{R2.1}{The ablation does not isolate the effect of component X.}
\resp{It does, but indirectly: rows 3 vs.\ 5 of Table~4 differ only in X.
We will reorder the table so the pair is adjacent and label it
``X on/off''.}

\rpoint{R2.3}{Novelty over [17] is unclear.}
\resp{[17] solves the offline variant; our setting is streaming with
deletions (Sec.~2, L96--104), which [17] explicitly leaves open
(their Sec.~7). The techniques also differ: we replace their global
re-indexing with the incremental structure of Sec.~3.2.}

\rhead{Reviewer 3 (R3)}

\rpoint{R3.1}{Could the method extend to setting Y?}
\resp{We believe so via the reduction in Appendix~D, but we have not run
this experiment and therefore make no empirical claim; we will note it as
future work.}

% --- closing line (optional if space allows) ----------------------------------
\vspace{4pt}
\noindent We will incorporate all clarifications above in the revision and
believe they are presentation-level changes within the camera-ready budget.

\end{document}
